% !TeX root = ./rapportUTT.tex
% !BIB TS-program = biber

\documentclass{rUTT}
% Pour retirer le thème couleur UTT,
%   Commenter la ligne précédente
%   Décommenter la ligne dessous
% \documentclass[noUTTcolors]{rUTT}


%%%%%%%%%%%%%%%%%%%%%%%%%%%%%%%%%%%%%%%%%%%%%%%%%%%%%%%%%%%%
% Quelques trucs à savoir pour modifier en paix :
% TOC = Table of Contents
% LOF / LOT = List Of Figures / List Of Tables
%%%%%%%%%%%%%%%%%%%%%%%%%%%%%%%%%%%%%%%%%%%%%%%%%%%%%%%%%%%%git@github.com:n3rada/ScribUTT.git
% Si on est on mode rapport de stage :
% Pour la confidentialité

\RPeda{{\sc ELGALAI} Reza} % Responsable pédagogique

\Entreprise{Login Sécurité} % Nom de l'entreprise / organisme / institution
\Lieu{199 Bureaux de la Colline, 92210 Saint-Cloud}
\REntre{Jean-Christophe PILLET}

% Mots clés du Thésaurus ou mots clés pour les metadatas du pdf
\Kone{Rétro-ingénierie} % Nature de l'activité
\Ktwo{Analyse de logiciels malveillants} % Branche d'activité économique
\Kthree{Détection d'intrusion} % Domaines technologiques
\Kfourth{Security Operation Center (SOC)} % Application Physique directe
%%%%%%%%%%%%%%%%%%%%%%%%%%%%%%%%%%%%%%%%%%%%%%%%%%%%%%%%%%%%

\Semestre{2025 - 2026}
%\UE{LT01} %Nom de l'UE OU nom complet de la branche si en mode rapport de stage !
\UE{Expert Forensic et Cybersécurité}

% Le titre de votre rapport OU le résumé de votre stage si en mode stage
% \newcommand{\titletext}{Un rapport en \LaTeX \\ écrit avec amour}

\newcommand{\titletext}{

    Ce mémoire rend compte d'un travail mené au cours de mon alternance chez Login Sécurité,
prestataire de services de sécurité managés (MSSP) du groupe Constellation, où
j'ai occupé le poste d'analyste SOC/VOC.

L'étude n'était pas au programme. Je l'ai engagée en marge de mes activités, sur un
échantillon dont la classe méritait selon moi une caractérisation complète, et les
règles qui en sortent peuvent alimenter la base de détection du parc client si
l'équipe les valide.

Le travail s'est déroulé en trois temps :

\begin{itemize}
    \item analyse statique du binaire, pour en extraire les caractéristiques structurelles ;
    \item analyse dynamique de son comportement en environnement confiné ;
    \item traduction de ces observations en règles YARA et Sigma, puis en plugin Volatility.
\end{itemize}

Un SOC qui attend la signature de l'éditeur ne détecte que ce qui est déjà
répertorié. Entre la découverte d'une menace et sa prise en charge par les produits
du marché, il s'écoule des jours. Savoir caractériser un binaire soi-même, c'est les
récupérer.
}



%% En mode Année - Mois - Jour
%\date{\today} % Pour la date de compilation
\date{15/08/2026}

\author{{\sc GOHAUT} Renan
% \and
% {\sc Nom} Prénom
% \break
% {\sc Nom} Prénom
% \and
% {\sc Nom} Prénom
% \break
% {\sc Nom} Prénom
}

% Texte affiché sur le carré bleu
\newcommand{\jobposition}{}
%\title{Vous êtes probablement assez bon pour travailler dans cette entreprise pour laquelle vous pensez ne pas être assez bon.}

\title{Détection de malwares à l'ère des menaces génératives en environnement SOC}

%%%%%%% Pour les métadonnées
% Fonctionne, seulement il y a un warning à cause des "\\" dans le title
% A tester ! Pour voir les infos du pdf on fait pdfinfo sous linux
% https://ctan.gutenberg.eu.org/macros/latex/contrib/hyperref/doc/hyperref-doc.pdf#subsection.3.7
\hypersetup{
    pdfauthor = {\theauthor}, % Auteur : pdfauthor = {\theauthor}
    pdftitle = {\thetitle}, % Titre : pdftitle = {Titre du document}
    pdfkeywords = {\theKone, \theKtwo, \theKthree, \theKfourth}
}

%%%%%%%
% n'oubliez pas de changer le language principal dans rUTT.cls
% en options du package babel !

\begin{document}
    %%%% - Choix de la page de garde

    %\frontpagereports % Pour le modèle rapports de TDs / TPs / Projets
    \frontpageSTB % Pour le modèle rapports de Stages des Branches / Master
    %\frontpageSTC % Pour le modèle rapports de Stages des TC

    {
        % \selectlanguage{english}
        \myILB
    }

    % page blanche après page de garde pour impression recto verso

    \pagestyle{UTT} % Pour que notre document utilise ce style
    % Ici on organise nos parties
    \justifying % on justifie notre texte via ragged2e


    \pagenumbering{gobble} % on n'affiche pas les numéros de page
    \import{latex-files/}{remerciements.tex} % Toujours avant le sommaire !

    \clearpage

    %%%%%%%%%%%%%%%%%%%%%%
    %% Sommaire
    %%%%%%%%%%%%%%%%%%%%%%

    {
        \setcounter{tocdepth}{1}
        \renewcommand{\contentsname}{Sommaire}
        \tableofcontents
    }

    \clearpage

    {
        % Tableaux et figures
        \phantomsection

        \listoftables
        \addcontentsline{toc}{chapter}{\listtablename}

        \listoffigures
        \addcontentsline{toc}{chapter}{\listfigurename}

        \markleft{Tableaux et figures}
    }

    % \setglossarystyle{listhypergroup}
    % \printglossary[type=\acronymtype, title={Liste des acronymes}, toctitle={Liste des acronymes}]
    % \printglossary[title={Glossaire}, toctitle={Glossaire}]

    \clearpage

    \pagenumbering{arabic}

    \import{latex-files/}{introduction.tex}

    \clearpage

    \import{latex-files/}{one.tex}
    \import{latex-files/}{two.tex}

    \clearpage

    \import{latex-files/}{conclusion.tex}
    \label{LastPage}

    \clearpage

    % Bibliographie !
    {
        \pagenumbering{gobble} % On n'affiche pas les numéros de page
        \phantomsection % hyperlinks will target the correct page
        \markboth{Bibliographie}{}
        \raggedright % pour éviter certaines erreurs rares d'affichage
        \sloppy
        \nocite{*} % pour faire apparaître tout du fichier bib
        \printbibliography[title={Bibliographie},heading=bibintoc]
    }

    \clearpage

    %\tripleS
    \pagenumbering{Roman} % On numérote en romain pour les annexes

    % Annexes !
    \import{latex-files/annexes/}{annexes.tex}

    \clearpage
    % Toujours avoir la table des matières en dernier ! Ici figure tout, même les annexes
    % Commenter/supprimer pour enlever la table des matières
    \myfinaltoc
    % On laisse une page blanche à la fin pour l'impression, c'est plus joli
    \clearpage
    \myemptypage

\end{document}
